\chapter{Resultados}\label{cap:resultados}
La \cref{fig:bode} se dibuja directamente desde \texttt{datos/mediciones.csv}.\index{diagrama de Bode}

\begin{figure}[htbp]
  \centering
  \begin{tikzpicture}
    \begin{semilogxaxis}[width=0.9\textwidth, height=6cm, grid=both,
        xlabel={Frecuencia (\si{\hertz})}, ylabel={$|H|$ (\si{\decibel})},
        legend pos=south west]
      \addplot[blue, thick] table[x=f, y=db, col sep=semicolon] {datos/mediciones.csv};
      \addlegendentry{medido}
      \addplot[red, dashed, domain=100:200000, samples=200]
        {20*log10(1/sqrt((1-(2*pi*x)^2*1e-9)^2 + (2*pi*x*47*1e-7)^2))};
      \addlegendentry{teórico}
      \draw[gray, dotted] (axis cs:5033,-60) -- (axis cs:5033,20) node[above, black, font=\small] {$f_0$};
    \end{semilogxaxis}
  \end{tikzpicture}
  \caption{Diagrama de Bode: medición (archivo \texttt{.csv}) frente a teoría.}\label{fig:bode}
\end{figure}

\section{Discusión}
El error de $f_0$ es \SI{0.45}{\percent} (\cref{tab:resumen}), dentro del objetivo. La forma de la
curva de la \cref{fig:senal} coincide con la respuesta subamortiguada de \textcite{alexander2013}.
\index{error de medición}\index[autores]{Ogata, K.}
